% BEGIN VOL08-EXPANSION VIII20
\section{Supplementary worked examples}
\label{sec:viii20-pedagogy}

\begin{example}[A deformation retract does not change homology]\label{ex:viii20-ped-01}
If $A\subset X$ is a deformation retract, the inclusion $i:A\hookrightarrow X$ and retraction $r:X\to A$ satisfy $r i=\operatorname{id}_A$ and $i r\simeq\operatorname{id}_X$. Homotopy invariance gives $i_*r_*=\operatorname{id}$ and $r_*i_*=\operatorname{id}$, so $H_*(A)\cong H_*(X)$.
\end{example}

\begin{example}[Contractible spaces have point homology]\label{ex:viii20-ped-02}
If $X$ is contractible, $\operatorname{id}_X$ is homotopic to a constant map. Hence $X$ has the homology of a point:
\[
H_0(X)\cong\mathbb Z,\qquad H_n(X)=0\quad(n>0)
\]
when $X$ is nonempty and path connected.
\end{example}

\begin{example}[Homotopic maps act identically on homology]\label{ex:viii20-ped-03}
For $f,g:X\to Y$ with $f\simeq g$, the induced chain maps are chain homotopic. Therefore $f_*=g_*:H_n(X)\to H_n(Y)$ in every degree. This is the algebraic mechanism behind homotopy invariance.
\end{example}

\section{Graded supplementary exercises}
The following exercises progress from direct computation and construction to proof, hypothesis testing, application, and synthesis.

\subsection*{Standard computations and constructions}

\begin{exercise}[Constant map in positive degree]\label{exr:viii20-ped-01}
What map does a constant map induce on $H_n$ for $n>0$?
\end{exercise}
\begin{hint}
It factors through a point.
\end{hint}
\begin{solution}
It induces the zero map because $H_n(\{*\})=0$ for $n>0$.
\end{solution}

\begin{exercise}[Interval homology]\label{exr:viii20-ped-02}
Compute $H_*(I)$ using homotopy invariance.
\end{exercise}
\begin{hint}
$I$ contracts to a point.
\end{hint}
\begin{solution}
$H_0(I)\cong\mathbb Z$ and $H_n(I)=0$ for $n>0$.
\end{solution}

\begin{exercise}[Cylinder homology]\label{exr:viii20-ped-03}
Compute $H_*(X\times I)$ in terms of $H_*(X)$.
\end{exercise}
\begin{hint}
Project to $X$.
\end{hint}
\begin{solution}
$X\times I$ deformation retracts onto $X\times\{0\}$, so $H_*(X\times I)\cong H_*(X)$.
\end{solution}

\begin{exercise}[Punctured plane]\label{exr:viii20-ped-04}
Compute the homology of $\mathbb R^2-\{0\}$.
\end{exercise}
\begin{hint}
Radially deformation retract onto $S^1$.
\end{hint}
\begin{solution}
$H_0\cong\mathbb Z$, $H_1\cong\mathbb Z$, and all higher groups vanish.
\end{solution}

\begin{exercise}[Annulus]\label{exr:viii20-ped-05}
Compute the homology of an annulus.
\end{exercise}
\begin{hint}
Deformation retract onto its core circle.
\end{hint}
\begin{solution}
It has the homology of $S^1$: $\mathbb Z$ in degrees $0$ and $1$, zero otherwise.
\end{solution}

\subsection*{Proofs}

\begin{exercise}[Homotopic maps]\label{exr:viii20-ped-06}
Prove that if $f\simeq g$, then $f_*=g_*$ on homology.
\end{exercise}
\begin{hint}
Use the prism operator or the chain-homotopy theorem.
\end{hint}
\begin{solution}
The homotopy produces a chain homotopy $P$ with $g_\#-f_\#=\partial P+P\partial$; boundaries differ by boundaries on cycles, so induced homology maps agree.
\end{solution}

\begin{exercise}[Homotopy equivalence]\label{exr:viii20-ped-07}
Prove a homotopy equivalence induces homology isomorphisms.
\end{exercise}
\begin{hint}
Apply homotopy invariance to the two homotopy-inverse composites.
\end{hint}
\begin{solution}
If $g$ is a homotopy inverse to $f$, then $g_*f_*=\operatorname{id}$ and $f_*g_*=\operatorname{id}$, so $f_*$ is an isomorphism.
\end{solution}

\begin{exercise}[Retract injection]\label{exr:viii20-ped-08}
If $A$ is a retract of $X$, prove $H_n(A)\to H_n(X)$ is injective.
\end{exercise}
\begin{hint}
Use $r\circ i=\operatorname{id}_A$.
\end{hint}
\begin{solution}
On homology, $r_*i_*=\operatorname{id}$, so $i_*$ has a left inverse and is injective.
\end{solution}

\begin{exercise}[Contractible reduced homology]\label{exr:viii20-ped-09}
Prove a nonempty contractible space has zero reduced homology.
\end{exercise}
\begin{hint}
Compare it with a point.
\end{hint}
\begin{solution}
Homotopy equivalence with a point gives $\widetilde H_n(X)=0$ in every degree.
\end{solution}

\subsection*{Counterexamples and hypothesis tests}

\begin{exercise}[Converse failure]\label{exr:viii20-ped-10}
Test: spaces with isomorphic homology groups are homotopy equivalent.
\end{exercise}
\begin{hint}
Homology is not a complete invariant.
\end{hint}
\begin{solution}
False. Distinct homotopy types can share all homology groups; homology forgets information such as much of the fundamental group and higher operations.
\end{solution}

\begin{exercise}[Acyclic means contractible]\label{exr:viii20-ped-11}
Test: if $\widetilde H_*(X)=0$, then $X$ is contractible.
\end{exercise}
\begin{hint}
Acyclicity is weaker than contractibility.
\end{hint}
\begin{solution}
False in general; there are noncontractible acyclic spaces.
\end{solution}

\begin{exercise}[Homology detects all homotopies]\label{exr:viii20-ped-12}
Test: if $f_*=g_*$ on all homology groups, then $f\simeq g$.
\end{exercise}
\begin{hint}
Induced homology maps are coarse invariants.
\end{hint}
\begin{solution}
False. Different homotopy classes can induce the same homology maps.
\end{solution}

\subsection*{Applications and investigations}

\begin{exercise}[Simplifying a space]\label{exr:viii20-ped-13}
Why is deformation retraction a computational tool rather than only a qualitative statement?
\end{exercise}
\begin{hint}
It replaces a space by a smaller homotopy-equivalent model.
\end{hint}
\begin{solution}
Any homology computation may be carried out on the retract, often reducing a complicated geometric space to a graph, sphere, or CW skeleton.
\end{solution}

\begin{exercise}[Model validation]\label{exr:viii20-ped-14}
How can homology invariance help check a numerical or combinatorial simplification of a shape?
\end{exercise}
\begin{hint}
A claimed homotopy-preserving simplification should preserve homology.
\end{hint}
\begin{solution}
If the simplification is intended to be a deformation retract or homotopy equivalence, differing Betti numbers expose a topological error immediately.
\end{solution}

\subsection*{Challenge problems}

\begin{exercise}[Wedge computation]\label{exr:viii20-ped-15}
Use deformation and known homology to determine $H_1$ of a figure-eight graph.
\end{exercise}
\begin{hint}
It is already a wedge $S^1\vee S^1$.
\end{hint}
\begin{solution}
$H_1\cong\mathbb Z^2$; the two loop classes are independent.
\end{solution}

\begin{exercise}[Invariant hierarchy]\label{exr:viii20-ped-16}
Explain why homotopy invariance is essential when defining any topological invariant from a chain model.
\end{exercise}
\begin{hint}
Different models of the same homotopy type should agree.
\end{hint}
\begin{solution}
It guarantees the chain-level construction descends to the homotopy category: homotopy-equivalent spaces receive canonically isomorphic homology groups.
\end{solution}

% END VOL08-EXPANSION VIII20
